%%%%%%%%%%%%%%%%%%%%%%%%%%%%%%%%%%%%%%%%%%%%%%%%%%%%%%%%%%%%
% 8.4 STATISTICAL ESTIMATION
% The Composition of Advanced Mathematics
%%%%%%%%%%%%%%%%%%%%%%%%%%%%%%%%%%%%%%%%%%%%%%%%%%%%%%%%%%%%

\section{Statistical Estimation}
\label{sec:statistical-estimation}

\prerequisite{
Descriptive statistics, sampling distributions, probability
distributions, expectation, variance, likelihood, derivatives,
optimization, laws of large numbers, central limit theorems, and basic
multivariable calculus.
}

\learningobjectives{
By the end of this section, the reader should be able to:
\begin{itemize}
    \item distinguish parameters, estimators, and estimates;
    \item define point estimation;
    \item compute bias, variance, and mean squared error;
    \item distinguish unbiasedness from consistency;
    \item understand efficiency and relative efficiency;
    \item understand sufficient statistics;
    \item use the factorization theorem;
    \item understand complete statistics conceptually;
    \item apply the Rao--Blackwell theorem;
    \item understand the Lehmann--Scheff\'e theorem;
    \item construct method-of-moments estimators;
    \item define likelihood and log-likelihood functions;
    \item compute maximum likelihood estimators;
    \item understand invariance of maximum likelihood estimators;
    \item derive score functions;
    \item define Fisher information;
    \item understand the Cram\'er--Rao lower bound;
    \item compare estimator efficiency;
    \item understand asymptotic consistency of maximum likelihood;
    \item understand asymptotic normality of estimators;
    \item compute approximate standard errors using Fisher information;
    \item understand observed information;
    \item recognize multiparameter estimation;
    \item distinguish scalar and matrix Fisher information;
    \item understand constrained estimation;
    \item recognize Bayesian estimators conceptually;
    \item understand shrinkage and regularization conceptually;
    \item recognize robust estimation;
    \item understand M-estimators conceptually;
    \item connect estimation with confidence intervals, hypothesis
          testing, regression, Bayesian statistics, and statistical
          learning.
\end{itemize}
}

%%%%%%%%%%%%%%%%%%%%%%%%%%%%%%%%%%%%%%%%%%%%%%%%%%%%%%%%%%%%
\subsection{The Goal of Statistical Estimation}
\label{subsec:goal-statistical-estimation}
%%%%%%%%%%%%%%%%%%%%%%%%%%%%%%%%%%%%%%%%%%%%%%%%%%%%%%%%%%%%

A statistical model contains unknown quantities called parameters.

For example, suppose

\[
X_1,\ldots,X_n
\overset{\text{i.i.d.}}{\sim}
N(\mu,\sigma^2).
\]

The quantities

\[
\mu
\]

and

\[
\sigma^2
\]

are population parameters.

Because these values are unknown, they must be estimated from the
observed sample.

Statistical estimation develops systematic methods for constructing and
evaluating such estimates.

%%%%%%%%%%%%%%%%%%%%%%%%%%%%%%%%%%%%%%%%%%%%%%%%%%%%%%%%%%%%
\subsection{Parameter Space}
\label{subsec:parameter-space}
%%%%%%%%%%%%%%%%%%%%%%%%%%%%%%%%%%%%%%%%%%%%%%%%%%%%%%%%%%%%

A statistical model may be written

\[
\boxed{
\{f(x;\theta):\theta\in\Theta\},
}
\]

where

\[
\theta
\]

is the parameter and

\[
\Theta
\]

is the parameter space.

For example, in a Bernoulli model,

\[
X\sim\operatorname{Bernoulli}(p),
\]

the parameter space is

\[
\boxed{
\Theta=[0,1].
}
\]

For a normal mean with known variance,

\[
\boxed{
\Theta=\mathbb{R}.
}
\]

%%%%%%%%%%%%%%%%%%%%%%%%%%%%%%%%%%%%%%%%%%%%%%%%%%%%%%%%%%%%
\subsection{Estimator and Estimate}
\label{subsec:estimator-estimate-statistical-estimation}
%%%%%%%%%%%%%%%%%%%%%%%%%%%%%%%%%%%%%%%%%%%%%%%%%%%%%%%%%%%%

\begin{definitionbox}{Estimator}
An estimator is a statistic

\[
\boxed{
\widehat{\theta}
=
T(X_1,\ldots,X_n)
}
\]

used to estimate an unknown parameter \(\theta\).
\end{definitionbox}

Before the sample is observed,

\[
\widehat{\theta}
\]

is a random variable.

After observing data

\[
x_1,\ldots,x_n,
\]

the value

\[
\widehat{\theta}(x_1,\ldots,x_n)
\]

is called an estimate.

%%%%%%%%%%%%%%%%%%%%%%%%%%%%%%%%%%%%%%%%%%%%%%%%%%%%%%%%%%%%
\subsection{Point Estimation}
\label{subsec:point-estimation}
%%%%%%%%%%%%%%%%%%%%%%%%%%%%%%%%%%%%%%%%%%%%%%%%%%%%%%%%%%%%

\begin{definitionbox}{Point Estimate}
A point estimate is a single numerical value used to estimate an unknown
parameter.
\end{definitionbox}

Examples include

\[
\boxed{
\widehat{\mu}
=
\overline{X},
}
\]

\[
\boxed{
\widehat{p}
=
\frac{X}{n},
}
\]

and

\[
\boxed{
\widehat{\sigma}^2
=
S^2.
}
\]

Point estimation gives one best-guess value rather than a range of
plausible values.

Confidence intervals are developed in the next section.

%%%%%%%%%%%%%%%%%%%%%%%%%%%%%%%%%%%%%%%%%%%%%%%%%%%%%%%%%%%%
\subsection{Properties of Good Estimators}
\label{subsec:properties-good-estimators}
%%%%%%%%%%%%%%%%%%%%%%%%%%%%%%%%%%%%%%%%%%%%%%%%%%%%%%%%%%%%

Several properties are commonly used to evaluate estimators:

\[
\boxed{
\text{unbiasedness},
}
\]

\[
\boxed{
\text{small variance},
}
\]

\[
\boxed{
\text{small mean squared error},
}
\]

\[
\boxed{
\text{consistency},
}
\]

\[
\boxed{
\text{efficiency},
}
\]

and

\[
\boxed{
\text{robustness}.
}
\]

No single property is universally sufficient.

%%%%%%%%%%%%%%%%%%%%%%%%%%%%%%%%%%%%%%%%%%%%%%%%%%%%%%%%%%%%
\subsection{Bias}
\label{subsec:bias-estimation}
%%%%%%%%%%%%%%%%%%%%%%%%%%%%%%%%%%%%%%%%%%%%%%%%%%%%%%%%%%%%

\begin{definitionbox}{Bias}
The bias of an estimator \(\widehat{\theta}\) is

\[
\boxed{
\operatorname{Bias}(\widehat{\theta})
=
\mathbb{E}[\widehat{\theta}]
-
\theta.
}
\]
\end{definitionbox}

If

\[
\mathbb{E}[\widehat{\theta}]
=
\theta,
\]

then \(\widehat{\theta}\) is unbiased.

%%%%%%%%%%%%%%%%%%%%%%%%%%%%%%%%%%%%%%%%%%%%%%%%%%%%%%%%%%%%
\subsection{Worked Example: Unbiased Sample Mean}
\label{subsec:worked-unbiased-sample-mean}
%%%%%%%%%%%%%%%%%%%%%%%%%%%%%%%%%%%%%%%%%%%%%%%%%%%%%%%%%%%%

\begin{workedexamplebox}{Sample Mean}

Suppose

\[
X_1,\ldots,X_n
\]

are i.i.d. with

\[
\mathbb{E}[X_i]=\mu.
\]

Then

\[
\overline{X}
=
\frac1n
\sum_{i=1}^{n}X_i.
\]

Taking expectations,

\[
\mathbb{E}[\overline{X}]
=
\frac1n
\sum_{i=1}^{n}\mu
=
\mu.
\]

\begin{answerbox}
\[
\boxed{
\operatorname{Bias}(\overline{X})=0.
}
\]
\end{answerbox}

Therefore the sample mean is unbiased for \(\mu\).

\end{workedexamplebox}

%%%%%%%%%%%%%%%%%%%%%%%%%%%%%%%%%%%%%%%%%%%%%%%%%%%%%%%%%%%%
\subsection{Biased Variance Estimator}
\label{subsec:biased-variance-estimator}
%%%%%%%%%%%%%%%%%%%%%%%%%%%%%%%%%%%%%%%%%%%%%%%%%%%%%%%%%%%%

Consider

\[
\widetilde{S}^2
=
\frac1n
\sum_{i=1}^{n}
(X_i-\overline{X})^2.
\]

Under i.i.d. sampling with finite variance,

\[
\boxed{
\mathbb{E}[\widetilde{S}^2]
=
\frac{n-1}{n}\sigma^2.
}
\]

Therefore,

\[
\boxed{
\operatorname{Bias}(\widetilde{S}^2)
=
-\frac{\sigma^2}{n}.
}
\]

This explains why division by \(n-1\) is used for the usual unbiased
sample variance.

%%%%%%%%%%%%%%%%%%%%%%%%%%%%%%%%%%%%%%%%%%%%%%%%%%%%%%%%%%%%
\subsection{Variance of an Estimator}
\label{subsec:variance-estimator-estimation}
%%%%%%%%%%%%%%%%%%%%%%%%%%%%%%%%%%%%%%%%%%%%%%%%%%%%%%%%%%%%

The variance

\[
\boxed{
\operatorname{Var}(\widehat{\theta})
}
\]

measures how much the estimator varies across repeated samples.

An estimator with smaller variance is more precise.

For example,

\[
\boxed{
\operatorname{Var}(\overline{X})
=
\frac{\sigma^2}{n}.
}
\]

Thus increasing sample size generally improves precision.

%%%%%%%%%%%%%%%%%%%%%%%%%%%%%%%%%%%%%%%%%%%%%%%%%%%%%%%%%%%%
\subsection{Mean Squared Error}
\label{subsec:mse-estimation}
%%%%%%%%%%%%%%%%%%%%%%%%%%%%%%%%%%%%%%%%%%%%%%%%%%%%%%%%%%%%

\begin{definitionbox}{Mean Squared Error}
The mean squared error of an estimator \(\widehat{\theta}\) is

\[
\boxed{
\operatorname{MSE}(\widehat{\theta})
=
\mathbb{E}
\left[
(\widehat{\theta}-\theta)^2
\right].
}
\]
\end{definitionbox}

The MSE combines bias and variance.

%%%%%%%%%%%%%%%%%%%%%%%%%%%%%%%%%%%%%%%%%%%%%%%%%%%%%%%%%%%%
\subsection{Bias--Variance Decomposition}
\label{subsec:bias-variance-decomposition}
%%%%%%%%%%%%%%%%%%%%%%%%%%%%%%%%%%%%%%%%%%%%%%%%%%%%%%%%%%%%

A fundamental identity is

\[
\boxed{
\operatorname{MSE}(\widehat{\theta})
=
\operatorname{Var}(\widehat{\theta})
+
\operatorname{Bias}(\widehat{\theta})^2.
}
\]

This means a small amount of bias may sometimes be acceptable if it
produces a large reduction in variance.

This idea becomes important in shrinkage and regularization.

%%%%%%%%%%%%%%%%%%%%%%%%%%%%%%%%%%%%%%%%%%%%%%%%%%%%%%%%%%%%
\subsection{Worked Example: Comparing Two Estimators}
\label{subsec:worked-comparing-estimators}
%%%%%%%%%%%%%%%%%%%%%%%%%%%%%%%%%%%%%%%%%%%%%%%%%%%%%%%%%%%%

\begin{workedexamplebox}{Bias Versus Variance}

Suppose estimator \(T_1\) has

\[
\operatorname{Bias}(T_1)=0
\]

and

\[
\operatorname{Var}(T_1)=9.
\]

Estimator \(T_2\) has

\[
\operatorname{Bias}(T_2)=1
\]

and

\[
\operatorname{Var}(T_2)=4.
\]

Then

\[
\operatorname{MSE}(T_1)=9,
\]

while

\[
\operatorname{MSE}(T_2)
=
4+1^2
=
5.
\]

\begin{answerbox}
\[
\boxed{
T_2
\text{ has smaller MSE despite being biased.}
}
\]
\end{answerbox}

\end{workedexamplebox}

%%%%%%%%%%%%%%%%%%%%%%%%%%%%%%%%%%%%%%%%%%%%%%%%%%%%%%%%%%%%
\subsection{Consistency}
\label{subsec:consistency-estimation}
%%%%%%%%%%%%%%%%%%%%%%%%%%%%%%%%%%%%%%%%%%%%%%%%%%%%%%%%%%%%

\begin{definitionbox}{Consistent Estimator}
An estimator sequence

\[
\widehat{\theta}_n
\]

is consistent for \(\theta\) if

\[
\boxed{
\widehat{\theta}_n
\xrightarrow{P}
\theta
}
\]

as

\[
n\to\infty.
\]
\end{definitionbox}

Consistency means the estimator becomes increasingly concentrated near
the true parameter as sample size grows.

%%%%%%%%%%%%%%%%%%%%%%%%%%%%%%%%%%%%%%%%%%%%%%%%%%%%%%%%%%%%
\subsection{Consistency of the Sample Mean}
\label{subsec:consistency-sample-mean}
%%%%%%%%%%%%%%%%%%%%%%%%%%%%%%%%%%%%%%%%%%%%%%%%%%%%%%%%%%%%

Under i.i.d. sampling with

\[
\mathbb{E}[|X_1|]<\infty,
\]

the law of large numbers gives

\[
\boxed{
\overline{X}_n
\xrightarrow{P}
\mu.
}
\]

Under the strong law,

\[
\boxed{
\overline{X}_n
\xrightarrow{\text{a.s.}}
\mu.
}
\]

Thus the sample mean is consistent for the population mean.

%%%%%%%%%%%%%%%%%%%%%%%%%%%%%%%%%%%%%%%%%%%%%%%%%%%%%%%%%%%%
\subsection{Mean-Square Consistency}
\label{subsec:mean-square-consistency}
%%%%%%%%%%%%%%%%%%%%%%%%%%%%%%%%%%%%%%%%%%%%%%%%%%%%%%%%%%%%

If

\[
\mathbb{E}
[
(\widehat{\theta}_n-\theta)^2
]
\to0,
\]

then

\[
\widehat{\theta}_n
\]

converges to \(\theta\) in mean square.

Since mean-square convergence implies convergence in probability,

\[
\boxed{
\operatorname{MSE}(\widehat{\theta}_n)\to0
\Longrightarrow
\widehat{\theta}_n\xrightarrow{P}\theta.
}
\]

%%%%%%%%%%%%%%%%%%%%%%%%%%%%%%%%%%%%%%%%%%%%%%%%%%%%%%%%%%%%
\subsection{Unbiasedness Does Not Imply Consistency}
\label{subsec:unbiased-not-consistent}
%%%%%%%%%%%%%%%%%%%%%%%%%%%%%%%%%%%%%%%%%%%%%%%%%%%%%%%%%%%%

An estimator may be unbiased for every \(n\) but fail to become more
concentrated around the parameter.

For consistency, variance must typically shrink appropriately.

Thus

\[
\boxed{
\text{unbiasedness}
\neq
\text{consistency}.
}
\]

%%%%%%%%%%%%%%%%%%%%%%%%%%%%%%%%%%%%%%%%%%%%%%%%%%%%%%%%%%%%
\subsection{Consistency Does Not Require Unbiasedness}
\label{subsec:consistent-biased}
%%%%%%%%%%%%%%%%%%%%%%%%%%%%%%%%%%%%%%%%%%%%%%%%%%%%%%%%%%%%

A sequence of estimators may have finite-sample bias that approaches zero.

If

\[
\operatorname{Bias}(\widehat{\theta}_n)\to0
\]

and

\[
\operatorname{Var}(\widehat{\theta}_n)\to0,
\]

then

\[
\operatorname{MSE}(\widehat{\theta}_n)\to0,
\]

which implies consistency.

%%%%%%%%%%%%%%%%%%%%%%%%%%%%%%%%%%%%%%%%%%%%%%%%%%%%%%%%%%%%
\subsection{Efficiency}
\label{subsec:estimator-efficiency}
%%%%%%%%%%%%%%%%%%%%%%%%%%%%%%%%%%%%%%%%%%%%%%%%%%%%%%%%%%%%

Suppose \(T_1\) and \(T_2\) are unbiased estimators of the same parameter.

If

\[
\operatorname{Var}(T_1)
<
\operatorname{Var}(T_2),
\]

then \(T_1\) is more efficient in the variance sense.

A relative efficiency measure may be written

\[
\boxed{
\operatorname{Eff}(T_1,T_2)
=
\frac{
\operatorname{Var}(T_2)
}{
\operatorname{Var}(T_1)
}.
}
\]

%%%%%%%%%%%%%%%%%%%%%%%%%%%%%%%%%%%%%%%%%%%%%%%%%%%%%%%%%%%%
\subsection{Sufficient Statistics}
\label{subsec:sufficient-statistics}
%%%%%%%%%%%%%%%%%%%%%%%%%%%%%%%%%%%%%%%%%%%%%%%%%%%%%%%%%%%%

\begin{definitionbox}{Sufficient Statistic}
A statistic \(T(X_1,\ldots,X_n)\) is sufficient for parameter \(\theta\)
if, conditional on \(T\), the distribution of the full sample does not
depend on \(\theta\).
\end{definitionbox}

Informally,

\[
\boxed{
T
\text{ contains all sample information about }\theta.
}
\]

Sufficiency provides a principled form of data reduction.

%%%%%%%%%%%%%%%%%%%%%%%%%%%%%%%%%%%%%%%%%%%%%%%%%%%%%%%%%%%%
\subsection{Factorization Theorem}
\label{subsec:factorization-theorem}
%%%%%%%%%%%%%%%%%%%%%%%%%%%%%%%%%%%%%%%%%%%%%%%%%%%%%%%%%%%%

\begin{theorem}[Neyman--Fisher Factorization Theorem]
A statistic \(T(X)\) is sufficient for \(\theta\) if the joint density or
mass function can be written as

\[
\boxed{
f(x_1,\ldots,x_n;\theta)
=
g(T(x_1,\ldots,x_n),\theta)
h(x_1,\ldots,x_n),
}
\]

where \(h\) does not depend on \(\theta\).
\end{theorem}

This theorem is one of the main tools for identifying sufficient
statistics.

%%%%%%%%%%%%%%%%%%%%%%%%%%%%%%%%%%%%%%%%%%%%%%%%%%%%%%%%%%%%
\subsection{Worked Example: Bernoulli Sufficiency}
\label{subsec:worked-bernoulli-sufficiency}
%%%%%%%%%%%%%%%%%%%%%%%%%%%%%%%%%%%%%%%%%%%%%%%%%%%%%%%%%%%%

\begin{workedexamplebox}{Bernoulli Sample}

Suppose

\[
X_1,\ldots,X_n
\overset{\text{i.i.d.}}{\sim}
\operatorname{Bernoulli}(p).
\]

The joint PMF is

\[
\prod_{i=1}^{n}
p^{x_i}(1-p)^{1-x_i}.
\]

Therefore,

\[
=
p^{\sum x_i}
(1-p)^{n-\sum x_i}.
\]

This depends on the sample through

\[
T
=
\sum_{i=1}^{n}X_i.
\]

By the factorization theorem,

\begin{answerbox}
\[
\boxed{
\sum_{i=1}^{n}X_i
\text{ is sufficient for }p.
}
\]
\end{answerbox}

\end{workedexamplebox}

%%%%%%%%%%%%%%%%%%%%%%%%%%%%%%%%%%%%%%%%%%%%%%%%%%%%%%%%%%%%
\subsection{Worked Example: Normal Mean Sufficiency}
\label{subsec:worked-normal-mean-sufficiency}
%%%%%%%%%%%%%%%%%%%%%%%%%%%%%%%%%%%%%%%%%%%%%%%%%%%%%%%%%%%%

\begin{workedexamplebox}{Normal Mean with Known Variance}

Suppose

\[
X_i
\overset{\text{i.i.d.}}{\sim}
N(\mu,\sigma^2),
\]

with known \(\sigma^2\).

The joint density is proportional to

\[
\exp
\left[
-\frac1{2\sigma^2}
\sum_{i=1}^{n}
(x_i-\mu)^2
\right].
\]

Expand:

\[
\sum(x_i-\mu)^2
=
\sum x_i^2
-
2\mu\sum x_i
+
n\mu^2.
\]

The dependence on \(\mu\) occurs through

\[
\sum x_i.
\]

Therefore,

\begin{answerbox}
\[
\boxed{
\sum_{i=1}^{n}X_i
\text{ and equivalently }\overline{X}
\text{ are sufficient for }\mu.
}
\]
\end{answerbox}

\end{workedexamplebox}

%%%%%%%%%%%%%%%%%%%%%%%%%%%%%%%%%%%%%%%%%%%%%%%%%%%%%%%%%%%%
\subsection{Minimal Sufficiency Preview}
\label{subsec:minimal-sufficiency-preview}
%%%%%%%%%%%%%%%%%%%%%%%%%%%%%%%%%%%%%%%%%%%%%%%%%%%%%%%%%%%%

A minimal sufficient statistic is a sufficient statistic that represents
the greatest possible reduction of the data without losing information
about the parameter.

Every other sufficient statistic must contain at least as much
parameter-relevant information in an appropriate sense.

Minimal sufficiency is developed more fully in mathematical statistics.

%%%%%%%%%%%%%%%%%%%%%%%%%%%%%%%%%%%%%%%%%%%%%%%%%%%%%%%%%%%%
\subsection{Complete Statistics}
\label{subsec:complete-statistics}
%%%%%%%%%%%%%%%%%%%%%%%%%%%%%%%%%%%%%%%%%%%%%%%%%%%%%%%%%%%%

\begin{definitionbox}{Complete Statistic}
A statistic \(T\) is complete for a family indexed by \(\theta\) if

\[
\mathbb{E}_\theta[g(T)]=0
\]

for every \(\theta\) implies

\[
\boxed{
P_\theta(g(T)=0)=1
}
\]

for every \(\theta\).
\end{definitionbox}

Completeness is a technical property that helps establish uniqueness of
optimal unbiased estimators.

%%%%%%%%%%%%%%%%%%%%%%%%%%%%%%%%%%%%%%%%%%%%%%%%%%%%%%%%%%%%
\subsection{Rao--Blackwell Theorem}
\label{subsec:rao-blackwell}
%%%%%%%%%%%%%%%%%%%%%%%%%%%%%%%%%%%%%%%%%%%%%%%%%%%%%%%%%%%%

\begin{theorem}[Rao--Blackwell Theorem]
Let \(U\) be an estimator of \(\theta\), and let \(T\) be sufficient for
\(\theta\).

Define

\[
\boxed{
U^*
=
\mathbb{E}[U\mid T].
}
\]

Then \(U^*\) has no larger mean squared error than \(U\) under squared
error loss.

If \(U\) is unbiased, then \(U^*\) is also unbiased and

\[
\boxed{
\operatorname{Var}(U^*)
\leq
\operatorname{Var}(U).
}
\]
\end{theorem}

Conditioning on a sufficient statistic can therefore improve an
estimator.

%%%%%%%%%%%%%%%%%%%%%%%%%%%%%%%%%%%%%%%%%%%%%%%%%%%%%%%%%%%%
\subsection{Why Rao--Blackwell Works}
\label{subsec:why-rao-blackwell}
%%%%%%%%%%%%%%%%%%%%%%%%%%%%%%%%%%%%%%%%%%%%%%%%%%%%%%%%%%%%

By the law of total variance,

\[
\operatorname{Var}(U)
=
\mathbb{E}
[
\operatorname{Var}(U\mid T)
]
+
\operatorname{Var}
(
\mathbb{E}[U\mid T]
).
\]

Therefore,

\[
\operatorname{Var}
(
\mathbb{E}[U\mid T]
)
\leq
\operatorname{Var}(U).
\]

Thus conditioning removes variation unrelated to the information in the
sufficient statistic.

%%%%%%%%%%%%%%%%%%%%%%%%%%%%%%%%%%%%%%%%%%%%%%%%%%%%%%%%%%%%
\subsection{Lehmann--Scheff\'e Theorem}
\label{subsec:lehmann-scheffe}
%%%%%%%%%%%%%%%%%%%%%%%%%%%%%%%%%%%%%%%%%%%%%%%%%%%%%%%%%%%%

\begin{theorem}[Lehmann--Scheff\'e Theorem]
If \(T\) is complete and sufficient for \(\theta\), and \(g(T)\) is an
unbiased estimator of a parameter function \(\tau(\theta)\), then
\(g(T)\) is the unique minimum-variance unbiased estimator of
\(\tau(\theta)\), up to almost-sure equality.
\end{theorem}

Such an estimator is commonly called a uniformly minimum-variance
unbiased estimator.

%%%%%%%%%%%%%%%%%%%%%%%%%%%%%%%%%%%%%%%%%%%%%%%%%%%%%%%%%%%%
\subsection{UMVU Estimators}
\label{subsec:umvu-estimators}
%%%%%%%%%%%%%%%%%%%%%%%%%%%%%%%%%%%%%%%%%%%%%%%%%%%%%%%%%%%%

UMVU stands for

\[
\boxed{
\text{Uniformly Minimum-Variance Unbiased}.
}
\]

An estimator \(\widehat{\tau}\) is UMVU if it is unbiased and has variance
no larger than any other unbiased estimator for every parameter value.

The Lehmann--Scheff\'e theorem provides a powerful route to proving this
property.

%%%%%%%%%%%%%%%%%%%%%%%%%%%%%%%%%%%%%%%%%%%%%%%%%%%%%%%%%%%%
\subsection{Method of Moments}
\label{subsec:method-of-moments}
%%%%%%%%%%%%%%%%%%%%%%%%%%%%%%%%%%%%%%%%%%%%%%%%%%%%%%%%%%%%

The method of moments estimates parameters by matching theoretical
moments to sample moments.

Suppose a model depends on parameters

\[
\theta_1,\ldots,\theta_k.
\]

Set

\[
\boxed{
\mathbb{E}_\theta[X^j]
=
\frac1n
\sum_{i=1}^{n}X_i^j
}
\]

for

\[
j=1,\ldots,k,
\]

and solve the resulting equations for the parameters.

%%%%%%%%%%%%%%%%%%%%%%%%%%%%%%%%%%%%%%%%%%%%%%%%%%%%%%%%%%%%
\subsection{Worked Example: Method of Moments for Poisson}
\label{subsec:worked-mom-poisson}
%%%%%%%%%%%%%%%%%%%%%%%%%%%%%%%%%%%%%%%%%%%%%%%%%%%%%%%%%%%%

\begin{workedexamplebox}{Poisson Rate}

Suppose

\[
X_1,\ldots,X_n
\overset{\text{i.i.d.}}{\sim}
\operatorname{Poisson}(\lambda).
\]

The theoretical mean is

\[
\mathbb{E}[X]=\lambda.
\]

Set this equal to the sample mean:

\[
\lambda
=
\overline{X}.
\]

Therefore the method-of-moments estimator is

\begin{answerbox}
\[
\boxed{
\widehat{\lambda}_{MM}
=
\overline{X}.
}
\]
\end{answerbox}

\end{workedexamplebox}

%%%%%%%%%%%%%%%%%%%%%%%%%%%%%%%%%%%%%%%%%%%%%%%%%%%%%%%%%%%%
\subsection{Worked Example: Method of Moments for Uniform}
\label{subsec:worked-mom-uniform}
%%%%%%%%%%%%%%%%%%%%%%%%%%%%%%%%%%%%%%%%%%%%%%%%%%%%%%%%%%%%

\begin{workedexamplebox}{Uniform Upper Endpoint}

Suppose

\[
X_i
\sim
\operatorname{Uniform}(0,\theta).
\]

Then

\[
\mathbb{E}[X]
=
\frac{\theta}{2}.
\]

Equating with the sample mean,

\[
\overline{X}
=
\frac{\theta}{2}.
\]

Thus,

\begin{answerbox}
\[
\boxed{
\widehat{\theta}_{MM}
=
2\overline{X}.
}
\]
\end{answerbox}

\end{workedexamplebox}

%%%%%%%%%%%%%%%%%%%%%%%%%%%%%%%%%%%%%%%%%%%%%%%%%%%%%%%%%%%%
\subsection{Likelihood Function}
\label{subsec:likelihood-function}
%%%%%%%%%%%%%%%%%%%%%%%%%%%%%%%%%%%%%%%%%%%%%%%%%%%%%%%%%%%%

Suppose observed data are

\[
x_1,\ldots,x_n.
\]

The likelihood function is the joint probability or density viewed as a
function of the unknown parameter.

\begin{definitionbox}{Likelihood}
For observed data \(x\), the likelihood is

\[
\boxed{
L(\theta;x)
=
f(x;\theta).
}
\]
\end{definitionbox}

The data are treated as fixed.

The parameter is treated as the variable.

%%%%%%%%%%%%%%%%%%%%%%%%%%%%%%%%%%%%%%%%%%%%%%%%%%%%%%%%%%%%
\subsection{Likelihood Is Not a Probability Distribution for the Parameter}
\label{subsec:likelihood-not-probability}
%%%%%%%%%%%%%%%%%%%%%%%%%%%%%%%%%%%%%%%%%%%%%%%%%%%%%%%%%%%%

In frequentist likelihood theory,

\[
L(\theta;x)
\]

is not generally normalized over \(\theta\).

Thus,

\[
\boxed{
L(\theta;x)
}
\]

does not represent

\[
P(\theta\mid x)
\]

in ordinary frequentist inference.

Posterior probabilities for parameters belong to Bayesian inference.

%%%%%%%%%%%%%%%%%%%%%%%%%%%%%%%%%%%%%%%%%%%%%%%%%%%%%%%%%%%%
\subsection{Log-Likelihood}
\label{subsec:log-likelihood}
%%%%%%%%%%%%%%%%%%%%%%%%%%%%%%%%%%%%%%%%%%%%%%%%%%%%%%%%%%%%

Because products can be difficult to differentiate, one usually works
with the logarithm of the likelihood:

\[
\boxed{
\ell(\theta)
=
\log L(\theta).
}
\]

Since logarithms are strictly increasing,

\[
\boxed{
\arg\max_\theta L(\theta)
=
\arg\max_\theta \ell(\theta).
}
\]

For independent observations, products become sums.

%%%%%%%%%%%%%%%%%%%%%%%%%%%%%%%%%%%%%%%%%%%%%%%%%%%%%%%%%%%%
\subsection{Maximum Likelihood Estimator}
\label{subsec:maximum-likelihood-estimator}
%%%%%%%%%%%%%%%%%%%%%%%%%%%%%%%%%%%%%%%%%%%%%%%%%%%%%%%%%%%%

\begin{definitionbox}{Maximum Likelihood Estimator}
The maximum likelihood estimator, abbreviated MLE, is a parameter value
that maximizes the likelihood:

\[
\boxed{
\widehat{\theta}_{MLE}
=
\arg\max_{\theta\in\Theta}
L(\theta).
}
\]
\end{definitionbox}

Equivalently,

\[
\boxed{
\widehat{\theta}_{MLE}
=
\arg\max_{\theta\in\Theta}
\ell(\theta).
}
\]

%%%%%%%%%%%%%%%%%%%%%%%%%%%%%%%%%%%%%%%%%%%%%%%%%%%%%%%%%%%%
\subsection{Maximum Likelihood Workflow}
\label{subsec:mle-workflow}
%%%%%%%%%%%%%%%%%%%%%%%%%%%%%%%%%%%%%%%%%%%%%%%%%%%%%%%%%%%%

A typical MLE procedure is:

\begin{enumerate}
    \item write the joint likelihood;
    \item simplify it;
    \item take the logarithm;
    \item differentiate with respect to the parameter;
    \item solve the score equation;
    \item check boundaries of the parameter space;
    \item verify that the solution gives a maximum.
\end{enumerate}

A derivative equal to zero identifies only a candidate point.

%%%%%%%%%%%%%%%%%%%%%%%%%%%%%%%%%%%%%%%%%%%%%%%%%%%%%%%%%%%%
\subsection{Worked Example: Bernoulli MLE}
\label{subsec:worked-bernoulli-mle}
%%%%%%%%%%%%%%%%%%%%%%%%%%%%%%%%%%%%%%%%%%%%%%%%%%%%%%%%%%%%

\begin{workedexamplebox}{Estimating a Bernoulli Probability}

Suppose

\[
X_1,\ldots,X_n
\overset{\text{i.i.d.}}{\sim}
\operatorname{Bernoulli}(p).
\]

The likelihood is

\[
L(p)
=
\prod_{i=1}^{n}
p^{x_i}
(1-p)^{1-x_i}.
\]

Let

\[
S
=
\sum_{i=1}^{n}x_i.
\]

Then

\[
L(p)
=
p^S(1-p)^{n-S}.
\]

The log-likelihood is

\[
\ell(p)
=
S\ln p
+
(n-S)\ln(1-p).
\]

Differentiate:

\[
\ell'(p)
=
\frac{S}{p}
-
\frac{n-S}{1-p}.
\]

Set equal to zero:

\[
\frac{S}{p}
=
\frac{n-S}{1-p}.
\]

Solving,

\[
S(1-p)=p(n-S),
\]

so

\[
S=np.
\]

Therefore,

\begin{answerbox}
\[
\boxed{
\widehat{p}_{MLE}
=
\frac Sn
=
\overline{X}.
}
\]
\end{answerbox}

\end{workedexamplebox}

%%%%%%%%%%%%%%%%%%%%%%%%%%%%%%%%%%%%%%%%%%%%%%%%%%%%%%%%%%%%
\subsection{Worked Example: Poisson MLE}
\label{subsec:worked-poisson-mle}
%%%%%%%%%%%%%%%%%%%%%%%%%%%%%%%%%%%%%%%%%%%%%%%%%%%%%%%%%%%%

\begin{workedexamplebox}{Estimating a Poisson Rate}

Suppose

\[
X_i
\overset{\text{i.i.d.}}{\sim}
\operatorname{Poisson}(\lambda).
\]

The likelihood is

\[
L(\lambda)
=
\prod_{i=1}^{n}
e^{-\lambda}
\frac{\lambda^{x_i}}{x_i!}.
\]

Thus,

\[
L(\lambda)
=
e^{-n\lambda}
\lambda^{\sum x_i}
\prod_{i=1}^{n}
\frac1{x_i!}.
\]

The log-likelihood is

\[
\ell(\lambda)
=
-n\lambda
+
\left(
\sum x_i
\right)
\ln\lambda
-
\sum\ln(x_i!).
\]

Differentiate:

\[
\ell'(\lambda)
=
-n
+
\frac{
\sum x_i
}{
\lambda
}.
\]

Set equal to zero:

\[
-n
+
\frac{
\sum x_i
}{
\lambda
}
=
0.
\]

Therefore,

\begin{answerbox}
\[
\boxed{
\widehat{\lambda}_{MLE}
=
\overline{X}.
}
\]
\end{answerbox}

\end{workedexamplebox}

%%%%%%%%%%%%%%%%%%%%%%%%%%%%%%%%%%%%%%%%%%%%%%%%%%%%%%%%%%%%
\subsection{Worked Example: Exponential MLE}
\label{subsec:worked-exponential-mle}
%%%%%%%%%%%%%%%%%%%%%%%%%%%%%%%%%%%%%%%%%%%%%%%%%%%%%%%%%%%%

\begin{workedexamplebox}{Exponential Rate Parameter}

Suppose

\[
X_i
\overset{\text{i.i.d.}}{\sim}
\operatorname{Exp}(\lambda),
\]

with rate parameter \(\lambda\).

The density is

\[
f(x;\lambda)
=
\lambda e^{-\lambda x},
\qquad
x\geq0.
\]

The likelihood is

\[
L(\lambda)
=
\lambda^n
e^{-\lambda\sum x_i}.
\]

The log-likelihood is

\[
\ell(\lambda)
=
n\ln\lambda
-
\lambda
\sum x_i.
\]

Differentiate:

\[
\ell'(\lambda)
=
\frac n\lambda
-
\sum x_i.
\]

Setting this equal to zero gives

\[
\lambda
=
\frac{
n
}{
\sum x_i
}.
\]

Thus,

\begin{answerbox}
\[
\boxed{
\widehat{\lambda}_{MLE}
=
\frac1{\overline{X}}.
}
\]
\end{answerbox}

\end{workedexamplebox}

%%%%%%%%%%%%%%%%%%%%%%%%%%%%%%%%%%%%%%%%%%%%%%%%%%%%%%%%%%%%
\subsection{Worked Example: Normal Mean MLE}
\label{subsec:worked-normal-mean-mle}
%%%%%%%%%%%%%%%%%%%%%%%%%%%%%%%%%%%%%%%%%%%%%%%%%%%%%%%%%%%%

\begin{workedexamplebox}{Normal Mean with Known Variance}

Suppose

\[
X_i
\overset{\text{i.i.d.}}{\sim}
N(\mu,\sigma^2),
\]

where \(\sigma^2\) is known.

Ignoring constants not involving \(\mu\),

\[
\ell(\mu)
=
-
\frac1{2\sigma^2}
\sum_{i=1}^{n}
(x_i-\mu)^2.
\]

Maximizing this is equivalent to minimizing

\[
\sum
(x_i-\mu)^2.
\]

Differentiate:

\[
\frac{d}{d\mu}
\sum
(x_i-\mu)^2
=
-2\sum(x_i-\mu).
\]

Setting equal to zero gives

\[
n\mu
=
\sum x_i.
\]

Therefore,

\begin{answerbox}
\[
\boxed{
\widehat{\mu}_{MLE}
=
\overline{X}.
}
\]
\end{answerbox}

\end{workedexamplebox}

%%%%%%%%%%%%%%%%%%%%%%%%%%%%%%%%%%%%%%%%%%%%%%%%%%%%%%%%%%%%
\subsection{Normal Variance MLE}
\label{subsec:normal-variance-mle}
%%%%%%%%%%%%%%%%%%%%%%%%%%%%%%%%%%%%%%%%%%%%%%%%%%%%%%%%%%%%

Suppose

\[
X_i
\overset{\text{i.i.d.}}{\sim}
N(\mu,\sigma^2),
\]

with both \(\mu\) and \(\sigma^2\) unknown.

The joint MLEs are

\[
\boxed{
\widehat{\mu}_{MLE}
=
\overline{X},
}
\]

and

\[
\boxed{
\widehat{\sigma}^2_{MLE}
=
\frac1n
\sum_{i=1}^{n}
(X_i-\overline{X})^2.
}
\]

The variance MLE uses denominator \(n\), not \(n-1\).

Therefore it is biased for finite \(n\).

%%%%%%%%%%%%%%%%%%%%%%%%%%%%%%%%%%%%%%%%%%%%%%%%%%%%%%%%%%%%
\subsection{MLE Need Not Be Unbiased}
\label{subsec:mle-not-unbiased}
%%%%%%%%%%%%%%%%%%%%%%%%%%%%%%%%%%%%%%%%%%%%%%%%%%%%%%%%%%%%

Maximum likelihood maximizes fit to the observed data.

It does not automatically guarantee unbiasedness.

For the normal variance,

\[
\mathbb{E}
[
\widehat{\sigma}^2_{MLE}
]
=
\frac{n-1}{n}\sigma^2.
\]

Thus,

\[
\boxed{
\widehat{\sigma}^2_{MLE}
}
\]

is biased downward.

Nevertheless, it is consistent.

%%%%%%%%%%%%%%%%%%%%%%%%%%%%%%%%%%%%%%%%%%%%%%%%%%%%%%%%%%%%
\subsection{Boundary Maximum Likelihood Estimates}
\label{subsec:boundary-mle}
%%%%%%%%%%%%%%%%%%%%%%%%%%%%%%%%%%%%%%%%%%%%%%%%%%%%%%%%%%%%

An MLE need not occur where

\[
\ell'(\theta)=0.
\]

The maximum may occur at a boundary.

For example, in a Bernoulli sample containing only failures,

\[
\sum X_i=0,
\]

the MLE is

\[
\boxed{
\widehat{p}=0.
}
\]

Similarly, if every observation is a success,

\[
\boxed{
\widehat{p}=1.
}
\]

Therefore parameter-space constraints must always be checked.

%%%%%%%%%%%%%%%%%%%%%%%%%%%%%%%%%%%%%%%%%%%%%%%%%%%%%%%%%%%%
\subsection{Uniform Distribution Endpoint MLE}
\label{subsec:uniform-endpoint-mle}
%%%%%%%%%%%%%%%%%%%%%%%%%%%%%%%%%%%%%%%%%%%%%%%%%%%%%%%%%%%%

Suppose

\[
X_i
\overset{\text{i.i.d.}}{\sim}
\operatorname{Uniform}(0,\theta).
\]

The density is

\[
f(x;\theta)
=
\frac1\theta
\mathbf{1}_{[0,\theta]}(x).
\]

The joint likelihood is

\[
L(\theta)
=
\theta^{-n}
\mathbf{1}_{\{
\theta\geq X_{(n)}
\}},
\]

where

\[
X_{(n)}
=
\max_iX_i.
\]

For allowable \(\theta\), the factor \(\theta^{-n}\) decreases as
\(\theta\) increases.

Therefore,

\[
\boxed{
\widehat{\theta}_{MLE}
=
X_{(n)}.
}
\]

%%%%%%%%%%%%%%%%%%%%%%%%%%%%%%%%%%%%%%%%%%%%%%%%%%%%%%%%%%%%
\subsection{Worked Example: MLE Versus Method of Moments}
\label{subsec:worked-mle-versus-mom}
%%%%%%%%%%%%%%%%%%%%%%%%%%%%%%%%%%%%%%%%%%%%%%%%%%%%%%%%%%%%

\begin{workedexamplebox}{Uniform Upper Endpoint}

For

\[
X_i
\sim
\operatorname{Uniform}(0,\theta),
\]

the method-of-moments estimator is

\[
\widehat{\theta}_{MM}
=
2\overline{X}.
\]

The maximum likelihood estimator is

\[
\widehat{\theta}_{MLE}
=
X_{(n)}.
\]

\begin{answerbox}
\[
\boxed{
\widehat{\theta}_{MM}=2\overline{X},
\qquad
\widehat{\theta}_{MLE}=X_{(n)}.
}
\]
\end{answerbox}

Different estimation principles can therefore produce different
estimators.

\end{workedexamplebox}

%%%%%%%%%%%%%%%%%%%%%%%%%%%%%%%%%%%%%%%%%%%%%%%%%%%%%%%%%%%%
\subsection{Invariance Property of MLEs}
\label{subsec:mle-invariance}
%%%%%%%%%%%%%%%%%%%%%%%%%%%%%%%%%%%%%%%%%%%%%%%%%%%%%%%%%%%%

Suppose

\[
\widehat{\theta}
\]

is the MLE of \(\theta\).

For a parameter transformation

\[
\eta=g(\theta),
\]

the MLE of \(\eta\) is

\[
\boxed{
\widehat{\eta}
=
g(\widehat{\theta})
}
\]

under the usual interpretation of the transformed parameter.

This is called the invariance property of maximum likelihood.

%%%%%%%%%%%%%%%%%%%%%%%%%%%%%%%%%%%%%%%%%%%%%%%%%%%%%%%%%%%%
\subsection{Worked Example: MLE Invariance}
\label{subsec:worked-mle-invariance}
%%%%%%%%%%%%%%%%%%%%%%%%%%%%%%%%%%%%%%%%%%%%%%%%%%%%%%%%%%%%

\begin{workedexamplebox}{Exponential Mean}

Suppose

\[
X_i
\sim
\operatorname{Exp}(\lambda),
\]

where \(\lambda\) is the rate.

The population mean is

\[
\mu
=
\frac1\lambda.
\]

Since

\[
\widehat{\lambda}_{MLE}
=
\frac1{\overline{X}},
\]

MLE invariance gives

\[
\widehat{\mu}_{MLE}
=
\frac1{
\widehat{\lambda}_{MLE}
}.
\]

Therefore,

\begin{answerbox}
\[
\boxed{
\widehat{\mu}_{MLE}
=
\overline{X}.
}
\]
\end{answerbox}

\end{workedexamplebox}

%%%%%%%%%%%%%%%%%%%%%%%%%%%%%%%%%%%%%%%%%%%%%%%%%%%%%%%%%%%%
\subsection{Score Function}
\label{subsec:score-function}
%%%%%%%%%%%%%%%%%%%%%%%%%%%%%%%%%%%%%%%%%%%%%%%%%%%%%%%%%%%%

\begin{definitionbox}{Score Function}
The score function is the derivative of the log-likelihood:

\[
\boxed{
U(\theta)
=
\frac{\partial}{\partial\theta}
\ell(\theta).
}
\]
\end{definitionbox}

For regular interior MLEs, the estimator often satisfies

\[
\boxed{
U(\widehat{\theta})=0.
}
\]

This is called the score equation.

%%%%%%%%%%%%%%%%%%%%%%%%%%%%%%%%%%%%%%%%%%%%%%%%%%%%%%%%%%%%
\subsection{Expected Score}
\label{subsec:expected-score}
%%%%%%%%%%%%%%%%%%%%%%%%%%%%%%%%%%%%%%%%%%%%%%%%%%%%%%%%%%%%

Under standard regularity conditions,

\[
\boxed{
\mathbb{E}_\theta[U(\theta)]
=
0.
}
\]

This follows by differentiating

\[
\int f(x;\theta)\,dx=1
\]

with respect to \(\theta\).

The result requires conditions allowing differentiation under the
integral sign.

%%%%%%%%%%%%%%%%%%%%%%%%%%%%%%%%%%%%%%%%%%%%%%%%%%%%%%%%%%%%
\subsection{Fisher Information}
\label{subsec:fisher-information}
%%%%%%%%%%%%%%%%%%%%%%%%%%%%%%%%%%%%%%%%%%%%%%%%%%%%%%%%%%%%

\begin{definitionbox}{Fisher Information}
For one observation, the Fisher information in parameter \(\theta\) is

\[
\boxed{
I(\theta)
=
\mathbb{E}_\theta
\left[
U(\theta)^2
\right].
}
\]
\end{definitionbox}

Under standard regularity conditions,

\[
\boxed{
I(\theta)
=
-
\mathbb{E}_\theta
\left[
\ell''(\theta)
\right].
}
\]

Fisher information measures how strongly the likelihood responds to
changes in the parameter.

%%%%%%%%%%%%%%%%%%%%%%%%%%%%%%%%%%%%%%%%%%%%%%%%%%%%%%%%%%%%
\subsection{Information in an Independent Sample}
\label{subsec:fisher-information-sample}
%%%%%%%%%%%%%%%%%%%%%%%%%%%%%%%%%%%%%%%%%%%%%%%%%%%%%%%%%%%%

For \(n\) independent identically distributed observations,

\[
\boxed{
I_n(\theta)
=
nI_1(\theta).
}
\]

Thus information grows linearly with sample size.

Consequently, standard estimation error typically shrinks at rate

\[
\boxed{
\frac1{\sqrt{n}}.
}
\]

%%%%%%%%%%%%%%%%%%%%%%%%%%%%%%%%%%%%%%%%%%%%%%%%%%%%%%%%%%%%
\subsection{Worked Example: Bernoulli Fisher Information}
\label{subsec:worked-bernoulli-information}
%%%%%%%%%%%%%%%%%%%%%%%%%%%%%%%%%%%%%%%%%%%%%%%%%%%%%%%%%%%%

\begin{workedexamplebox}{Fisher Information for Bernoulli Data}

For one Bernoulli observation,

\[
\ell(p)
=
X\ln p
+
(1-X)\ln(1-p).
\]

The score is

\[
U(p)
=
\frac{X}{p}
-
\frac{1-X}{1-p}.
\]

The second derivative is

\[
\ell''(p)
=
-\frac{X}{p^2}
-
\frac{1-X}{(1-p)^2}.
\]

Taking negative expectation,

\[
I(p)
=
\frac{p}{p^2}
+
\frac{1-p}{(1-p)^2}.
\]

Therefore,

\begin{answerbox}
\[
\boxed{
I(p)
=
\frac1{
p(1-p)
}.
}
\]
\end{answerbox}

For \(n\) observations,

\[
\boxed{
I_n(p)
=
\frac{n}{
p(1-p)
}.
}
\]

\end{workedexamplebox}

%%%%%%%%%%%%%%%%%%%%%%%%%%%%%%%%%%%%%%%%%%%%%%%%%%%%%%%%%%%%
\subsection{Worked Example: Poisson Fisher Information}
\label{subsec:worked-poisson-information}
%%%%%%%%%%%%%%%%%%%%%%%%%%%%%%%%%%%%%%%%%%%%%%%%%%%%%%%%%%%%

\begin{workedexamplebox}{Poisson Information}

For one Poisson observation,

\[
\ell(\lambda)
=
-\lambda
+
X\ln\lambda
-
\ln(X!).
\]

Then

\[
\ell''(\lambda)
=
-\frac{X}{\lambda^2}.
\]

Therefore,

\[
I(\lambda)
=
-
\mathbb{E}
[
\ell''(\lambda)
].
\]

Since

\[
\mathbb{E}[X]=\lambda,
\]

we obtain

\begin{answerbox}
\[
\boxed{
I(\lambda)
=
\frac1\lambda.
}
\]
\end{answerbox}

For \(n\) observations,

\[
\boxed{
I_n(\lambda)
=
\frac n\lambda.
}
\]

\end{workedexamplebox}

%%%%%%%%%%%%%%%%%%%%%%%%%%%%%%%%%%%%%%%%%%%%%%%%%%%%%%%%%%%%
\subsection{Cram\'er--Rao Lower Bound}
\label{subsec:cramer-rao}
%%%%%%%%%%%%%%%%%%%%%%%%%%%%%%%%%%%%%%%%%%%%%%%%%%%%%%%%%%%%

\begin{theorem}[Cram\'er--Rao Lower Bound]
Under regularity conditions, if \(\widehat{\theta}\) is an unbiased
estimator of a scalar parameter \(\theta\), then

\[
\boxed{
\operatorname{Var}(\widehat{\theta})
\geq
\frac1{I_n(\theta)}.
}
\]
\end{theorem}

More generally, for an unbiased estimator of \(g(\theta)\),

\[
\boxed{
\operatorname{Var}
(
\widehat{g(\theta)}
)
\geq
\frac{
[g'(\theta)]^2
}{
I_n(\theta)
}.
}
\]

%%%%%%%%%%%%%%%%%%%%%%%%%%%%%%%%%%%%%%%%%%%%%%%%%%%%%%%%%%%%
\subsection{Efficient Estimators}
\label{subsec:efficient-estimators-crlb}
%%%%%%%%%%%%%%%%%%%%%%%%%%%%%%%%%%%%%%%%%%%%%%%%%%%%%%%%%%%%

An unbiased estimator that achieves the Cram\'er--Rao lower bound is
called efficient in the Cram\'er--Rao sense.

Thus if

\[
\operatorname{Var}(\widehat{\theta})
=
\frac1{I_n(\theta)},
\]

then

\[
\boxed{
\widehat{\theta}
\text{ is efficient}.
}
\]

%%%%%%%%%%%%%%%%%%%%%%%%%%%%%%%%%%%%%%%%%%%%%%%%%%%%%%%%%%%%
\subsection{Worked Example: Efficiency of Bernoulli Sample Proportion}
\label{subsec:worked-bernoulli-efficiency}
%%%%%%%%%%%%%%%%%%%%%%%%%%%%%%%%%%%%%%%%%%%%%%%%%%%%%%%%%%%%

\begin{workedexamplebox}{Sample Proportion}

For Bernoulli data,

\[
I_n(p)
=
\frac{n}{
p(1-p)
}.
\]

Therefore the Cram\'er--Rao lower bound is

\[
\frac1{I_n(p)}
=
\frac{
p(1-p)
}{n}.
\]

But

\[
\operatorname{Var}(\widehat{p})
=
\frac{
p(1-p)
}{n}.
\]

Hence,

\begin{answerbox}
\[
\boxed{
\widehat{p}
=
\overline{X}
\text{ attains the Cram\'er--Rao bound}.
}
\]
\end{answerbox}

Thus it is efficient among unbiased estimators under the regular model.

\end{workedexamplebox}

%%%%%%%%%%%%%%%%%%%%%%%%%%%%%%%%%%%%%%%%%%%%%%%%%%%%%%%%%%%%
\subsection{Observed Information}
\label{subsec:observed-information}
%%%%%%%%%%%%%%%%%%%%%%%%%%%%%%%%%%%%%%%%%%%%%%%%%%%%%%%%%%%%

The observed information is

\[
\boxed{
J(\theta)
=
-
\ell''(\theta).
}
\]

Unlike Fisher information, which is an expectation under the model,
observed information depends on the realized data.

Evaluating at the MLE gives

\[
\boxed{
J(\widehat{\theta}).
}
\]

This is often used to approximate the variance of the MLE.

%%%%%%%%%%%%%%%%%%%%%%%%%%%%%%%%%%%%%%%%%%%%%%%%%%%%%%%%%%%%
\subsection{Asymptotic Normality of MLEs}
\label{subsec:asymptotic-normality-mle}
%%%%%%%%%%%%%%%%%%%%%%%%%%%%%%%%%%%%%%%%%%%%%%%%%%%%%%%%%%%%

Under suitable regularity conditions,

\[
\boxed{
\sqrt{n}
(
\widehat{\theta}_{MLE}
-
\theta
)
\xrightarrow{d}
N
\left(
0,
\frac1{I_1(\theta)}
\right).
}
\]

Equivalently,

\[
\boxed{
\widehat{\theta}_{MLE}
\approx
N
\left(
\theta,
\frac1{nI_1(\theta)}
\right).
}
\]

This is one of the most important asymptotic results in statistics.

%%%%%%%%%%%%%%%%%%%%%%%%%%%%%%%%%%%%%%%%%%%%%%%%%%%%%%%%%%%%
\subsection{Approximate Standard Error of the MLE}
\label{subsec:approx-se-mle}
%%%%%%%%%%%%%%%%%%%%%%%%%%%%%%%%%%%%%%%%%%%%%%%%%%%%%%%%%%%%

The asymptotic standard error is

\[
\boxed{
\operatorname{SE}
(
\widehat{\theta}_{MLE}
)
\approx
\frac1{
\sqrt{
I_n(\theta)
}
}.
}
\]

Since \(\theta\) is unknown, one may plug in the MLE:

\[
\boxed{
\widehat{\operatorname{SE}}
(
\widehat{\theta}_{MLE}
)
\approx
\frac1{
\sqrt{
I_n(\widehat{\theta})
}
}.
}
\]

Observed information may also be used.

%%%%%%%%%%%%%%%%%%%%%%%%%%%%%%%%%%%%%%%%%%%%%%%%%%%%%%%%%%%%
\subsection{Worked Example: Approximate SE of Bernoulli MLE}
\label{subsec:worked-se-bernoulli-mle}
%%%%%%%%%%%%%%%%%%%%%%%%%%%%%%%%%%%%%%%%%%%%%%%%%%%%%%%%%%%%

\begin{workedexamplebox}{Bernoulli Standard Error}

For Bernoulli data,

\[
I_n(p)
=
\frac{n}{
p(1-p)
}.
\]

Therefore,

\[
\operatorname{SE}(\widehat{p})
\approx
\sqrt{
\frac{
p(1-p)
}{n}
}.
\]

Replacing \(p\) by \(\widehat{p}\),

\begin{answerbox}
\[
\boxed{
\widehat{\operatorname{SE}}(\widehat{p})
=
\sqrt{
\frac{
\widehat{p}(1-\widehat{p})
}{n}
}.
}
\]
\end{answerbox}

\end{workedexamplebox}

%%%%%%%%%%%%%%%%%%%%%%%%%%%%%%%%%%%%%%%%%%%%%%%%%%%%%%%%%%%%
\subsection{Consistency of Maximum Likelihood}
\label{subsec:mle-consistency}
%%%%%%%%%%%%%%%%%%%%%%%%%%%%%%%%%%%%%%%%%%%%%%%%%%%%%%%%%%%%

Under suitable identifiability and regularity assumptions,

\[
\boxed{
\widehat{\theta}_{MLE}
\xrightarrow{P}
\theta.
}
\]

The intuition is that the average log-likelihood converges to its
population expectation, which is maximized at the true parameter under
correct model specification.

Consistency is asymptotic and does not imply perfect finite-sample
behavior.

%%%%%%%%%%%%%%%%%%%%%%%%%%%%%%%%%%%%%%%%%%%%%%%%%%%%%%%%%%%%
\subsection{Likelihood Curvature}
\label{subsec:likelihood-curvature}
%%%%%%%%%%%%%%%%%%%%%%%%%%%%%%%%%%%%%%%%%%%%%%%%%%%%%%%%%%%%

A sharply peaked log-likelihood indicates that small parameter changes
produce large decreases in fit.

A flat log-likelihood indicates more uncertainty.

The second derivative

\[
-\ell''(\theta)
\]

measures local curvature.

Thus greater information corresponds to greater likelihood curvature
near the true value.

%%%%%%%%%%%%%%%%%%%%%%%%%%%%%%%%%%%%%%%%%%%%%%%%%%%%%%%%%%%%
\subsection{Likelihood Ratio}
\label{subsec:likelihood-ratio-estimation}
%%%%%%%%%%%%%%%%%%%%%%%%%%%%%%%%%%%%%%%%%%%%%%%%%%%%%%%%%%%%

Likelihoods can also compare parameter values.

The relative likelihood of \(\theta\) compared with the MLE is

\[
\boxed{
R(\theta)
=
\frac{
L(\theta)
}{
L(\widehat{\theta})
}.
}
\]

Since the MLE maximizes the likelihood,

\[
\boxed{
0\leq R(\theta)\leq1.
}
\]

Values close to \(1\) indicate parameter values fitting nearly as well as
the MLE.

%%%%%%%%%%%%%%%%%%%%%%%%%%%%%%%%%%%%%%%%%%%%%%%%%%%%%%%%%%%%
\subsection{Profile Likelihood Preview}
\label{subsec:profile-likelihood-preview}
%%%%%%%%%%%%%%%%%%%%%%%%%%%%%%%%%%%%%%%%%%%%%%%%%%%%%%%%%%%%

Suppose the parameter vector is

\[
(\theta,\eta),
\]

where \(\theta\) is the parameter of interest and \(\eta\) is a nuisance
parameter.

For each fixed \(\theta\), maximize the likelihood over \(\eta\):

\[
\widehat{\eta}(\theta)
=
\arg\max_\eta
L(\theta,\eta).
\]

The profile likelihood is

\[
\boxed{
L_p(\theta)
=
L
(
\theta,
\widehat{\eta}(\theta)
).
}
\]

This reduces a multiparameter problem to a function of the parameter of
interest.

%%%%%%%%%%%%%%%%%%%%%%%%%%%%%%%%%%%%%%%%%%%%%%%%%%%%%%%%%%%%
\subsection{Nuisance Parameters}
\label{subsec:nuisance-parameters}
%%%%%%%%%%%%%%%%%%%%%%%%%%%%%%%%%%%%%%%%%%%%%%%%%%%%%%%%%%%%

A nuisance parameter affects the model but is not the main quantity of
scientific interest.

For example, in

\[
N(\mu,\sigma^2),
\]

one may care primarily about \(\mu\) while \(\sigma^2\) remains unknown.

Then \(\sigma^2\) acts as a nuisance parameter.

Statistical procedures often require estimating, profiling, conditioning,
or integrating over nuisance parameters.

%%%%%%%%%%%%%%%%%%%%%%%%%%%%%%%%%%%%%%%%%%%%%%%%%%%%%%%%%%%%
\subsection{Multiparameter Estimation}
\label{subsec:multiparameter-estimation}
%%%%%%%%%%%%%%%%%%%%%%%%%%%%%%%%%%%%%%%%%%%%%%%%%%%%%%%%%%%%

Suppose

\[
\boldsymbol{\theta}
=
(\theta_1,\ldots,\theta_k)^T.
\]

The score becomes a vector:

\[
\boxed{
\mathbf{U}(\boldsymbol{\theta})
=
\nabla_{\boldsymbol{\theta}}
\ell(\boldsymbol{\theta}).
}
\]

The MLE often satisfies

\[
\boxed{
\mathbf{U}
(
\widehat{\boldsymbol{\theta}}
)
=
\mathbf{0}.
}
\]

%%%%%%%%%%%%%%%%%%%%%%%%%%%%%%%%%%%%%%%%%%%%%%%%%%%%%%%%%%%%
\subsection{Fisher Information Matrix}
\label{subsec:fisher-information-matrix}
%%%%%%%%%%%%%%%%%%%%%%%%%%%%%%%%%%%%%%%%%%%%%%%%%%%%%%%%%%%%

For a parameter vector

\[
\boldsymbol{\theta},
\]

the Fisher information matrix is

\[
\boxed{
I(\boldsymbol{\theta})
=
\mathbb{E}
\left[
\mathbf{U}
\mathbf{U}^T
\right].
}
\]

Under regularity conditions,

\[
\boxed{
I(\boldsymbol{\theta})
=
-
\mathbb{E}
\left[
\nabla^2
\ell(\boldsymbol{\theta})
\right].
}
\]

The inverse information matrix approximates the covariance matrix of the
MLE.

%%%%%%%%%%%%%%%%%%%%%%%%%%%%%%%%%%%%%%%%%%%%%%%%%%%%%%%%%%%%
\subsection{Multivariate MLE Asymptotic Normality}
\label{subsec:multivariate-mle-normality}
%%%%%%%%%%%%%%%%%%%%%%%%%%%%%%%%%%%%%%%%%%%%%%%%%%%%%%%%%%%%

Under suitable regularity conditions,

\[
\boxed{
\sqrt{n}
\left(
\widehat{\boldsymbol{\theta}}
-
\boldsymbol{\theta}
\right)
\xrightarrow{d}
N_k
\left(
\mathbf{0},
I_1(\boldsymbol{\theta})^{-1}
\right).
}
\]

Thus approximately,

\[
\boxed{
\operatorname{Cov}
(
\widehat{\boldsymbol{\theta}}
)
\approx
I_n(\boldsymbol{\theta})^{-1}.
}
\]

%%%%%%%%%%%%%%%%%%%%%%%%%%%%%%%%%%%%%%%%%%%%%%%%%%%%%%%%%%%%
\subsection{Parameter Identifiability}
\label{subsec:parameter-identifiability}
%%%%%%%%%%%%%%%%%%%%%%%%%%%%%%%%%%%%%%%%%%%%%%%%%%%%%%%%%%%%

A model is identifiable if distinct parameter values correspond to
distinct probability distributions.

Formally,

\[
\boxed{
P_{\theta_1}
=
P_{\theta_2}
\Longrightarrow
\theta_1=\theta_2.
}
\]

If a model is not identifiable, data cannot uniquely distinguish certain
parameter values.

Identifiability is fundamental for consistent estimation.

%%%%%%%%%%%%%%%%%%%%%%%%%%%%%%%%%%%%%%%%%%%%%%%%%%%%%%%%%%%%
\subsection{Worked Example: Nonidentifiability}
\label{subsec:worked-nonidentifiability}
%%%%%%%%%%%%%%%%%%%%%%%%%%%%%%%%%%%%%%%%%%%%%%%%%%%%%%%%%%%%

\begin{workedexamplebox}{Squared Parameter}

Suppose a distribution depends only on

\[
\theta^2.
\]

Then parameters

\[
\theta
\]

and

\[
-\theta
\]

produce the same distribution.

Unless the parameter space restricts the sign,

\begin{answerbox}
\[
\boxed{
\theta
\text{ is not identifiable}.
}
\]
\end{answerbox}

The identifiable quantity is \(\theta^2\), or perhaps \(|\theta|\).

\end{workedexamplebox}

%%%%%%%%%%%%%%%%%%%%%%%%%%%%%%%%%%%%%%%%%%%%%%%%%%%%%%%%%%%%
\subsection{Constrained Estimation}
\label{subsec:constrained-estimation}
%%%%%%%%%%%%%%%%%%%%%%%%%%%%%%%%%%%%%%%%%%%%%%%%%%%%%%%%%%%%

Parameters often satisfy constraints.

Examples include

\[
\boxed{
p\in[0,1],
}
\]

\[
\boxed{
\sigma^2>0,
}
\]

and

\[
\boxed{
\sum_{i=1}^{k}p_i=1.
}
\]

Optimization must respect the parameter space.

An unconstrained critical point outside \(\Theta\) cannot be a valid
estimate.

%%%%%%%%%%%%%%%%%%%%%%%%%%%%%%%%%%%%%%%%%%%%%%%%%%%%%%%%%%%%
\subsection{Loss Functions}
\label{subsec:loss-functions-estimation}
%%%%%%%%%%%%%%%%%%%%%%%%%%%%%%%%%%%%%%%%%%%%%%%%%%%%%%%%%%%%

Estimation can be viewed as a decision problem.

A loss function measures the cost of estimating \(\theta\) by \(a\).

Squared-error loss is

\[
\boxed{
L(\theta,a)
=
(a-\theta)^2.
}
\]

Absolute-error loss is

\[
\boxed{
L(\theta,a)
=
|a-\theta|.
}
\]

Different loss functions can favor different estimators.

%%%%%%%%%%%%%%%%%%%%%%%%%%%%%%%%%%%%%%%%%%%%%%%%%%%%%%%%%%%%
\subsection{Risk Function}
\label{subsec:risk-function}
%%%%%%%%%%%%%%%%%%%%%%%%%%%%%%%%%%%%%%%%%%%%%%%%%%%%%%%%%%%%

The risk of an estimator is its expected loss:

\[
\boxed{
R(\theta,\widehat{\theta})
=
\mathbb{E}_\theta
[
L(\theta,\widehat{\theta})
].
}
\]

Under squared-error loss,

\[
\boxed{
R(\theta,\widehat{\theta})
=
\operatorname{MSE}(\widehat{\theta}).
}
\]

Decision-theoretic estimation compares estimators through their risk
functions.

%%%%%%%%%%%%%%%%%%%%%%%%%%%%%%%%%%%%%%%%%%%%%%%%%%%%%%%%%%%%
\subsection{Bayesian Estimation Preview}
\label{subsec:bayesian-estimation-preview}
%%%%%%%%%%%%%%%%%%%%%%%%%%%%%%%%%%%%%%%%%%%%%%%%%%%%%%%%%%%%

Bayesian estimation treats the parameter as uncertain before observing
the data.

A prior distribution

\[
\pi(\theta)
\]

is combined with the likelihood:

\[
L(\theta;x).
\]

Bayes' theorem gives the posterior

\[
\boxed{
\pi(\theta\mid x)
\propto
L(\theta;x)\pi(\theta).
}
\]

Point estimates can then be chosen from the posterior distribution.

%%%%%%%%%%%%%%%%%%%%%%%%%%%%%%%%%%%%%%%%%%%%%%%%%%%%%%%%%%%%
\subsection{Posterior Mean}
\label{subsec:posterior-mean-estimator}
%%%%%%%%%%%%%%%%%%%%%%%%%%%%%%%%%%%%%%%%%%%%%%%%%%%%%%%%%%%%

Under squared-error loss, the Bayes estimator is the posterior mean:

\[
\boxed{
\widehat{\theta}_{Bayes}
=
\mathbb{E}[\theta\mid X].
}
\]

Under absolute-error loss, a posterior median is optimal.

Under zero-one-type loss in suitable settings, a posterior mode may be
used.

Thus the preferred estimator depends on the loss function.

%%%%%%%%%%%%%%%%%%%%%%%%%%%%%%%%%%%%%%%%%%%%%%%%%%%%%%%%%%%%
\subsection{MAP Estimation}
\label{subsec:map-estimation}
%%%%%%%%%%%%%%%%%%%%%%%%%%%%%%%%%%%%%%%%%%%%%%%%%%%%%%%%%%%%

The maximum a posteriori estimate, abbreviated MAP, maximizes the
posterior density:

\[
\boxed{
\widehat{\theta}_{MAP}
=
\arg\max_\theta
\pi(\theta\mid x).
}
\]

Because

\[
\pi(\theta\mid x)
\propto
L(\theta;x)\pi(\theta),
\]

the MAP estimate balances likelihood and prior information.

If the prior is effectively flat over the relevant region, MAP and MLE
may coincide.

%%%%%%%%%%%%%%%%%%%%%%%%%%%%%%%%%%%%%%%%%%%%%%%%%%%%%%%%%%%%
\subsection{Shrinkage Estimation}
\label{subsec:shrinkage-estimation}
%%%%%%%%%%%%%%%%%%%%%%%%%%%%%%%%%%%%%%%%%%%%%%%%%%%%%%%%%%%%

A shrinkage estimator deliberately moves an estimate toward a target.

A simple form is

\[
\boxed{
\widehat{\theta}_{shrink}
=
(1-\lambda)
\widehat{\theta}
+
\lambda\theta_0,
}
\]

where

\[
0\leq\lambda\leq1.
\]

Shrinkage introduces bias but can reduce variance enough to lower mean
squared error.

%%%%%%%%%%%%%%%%%%%%%%%%%%%%%%%%%%%%%%%%%%%%%%%%%%%%%%%%%%%%
\subsection{Regularization}
\label{subsec:regularization-estimation}
%%%%%%%%%%%%%%%%%%%%%%%%%%%%%%%%%%%%%%%%%%%%%%%%%%%%%%%%%%%%

Regularized estimation modifies an objective function by penalizing
undesirable parameter values.

A generic form is

\[
\boxed{
\widehat{\theta}
=
\arg\min_\theta
\left\{
\text{data loss}
+
\lambda
\text{ penalty}
\right\}.
}
\]

Regularization is central to modern regression and statistical learning.

%%%%%%%%%%%%%%%%%%%%%%%%%%%%%%%%%%%%%%%%%%%%%%%%%%%%%%%%%%%%
\subsection{Ridge-Type Estimation Preview}
\label{subsec:ridge-preview-estimation}
%%%%%%%%%%%%%%%%%%%%%%%%%%%%%%%%%%%%%%%%%%%%%%%%%%%%%%%%%%%%

A quadratic penalty has the form

\[
\boxed{
\lambda
\|\boldsymbol{\beta}\|_2^2.
}
\]

This shrinks coefficient estimates toward zero.

It generally increases bias but can substantially reduce variance when
predictors are highly correlated or dimension is large.

%%%%%%%%%%%%%%%%%%%%%%%%%%%%%%%%%%%%%%%%%%%%%%%%%%%%%%%%%%%%
\subsection{Lasso-Type Estimation Preview}
\label{subsec:lasso-preview-estimation}
%%%%%%%%%%%%%%%%%%%%%%%%%%%%%%%%%%%%%%%%%%%%%%%%%%%%%%%%%%%%

An absolute-value penalty has the form

\[
\boxed{
\lambda
\|\boldsymbol{\beta}\|_1.
}
\]

This can produce some coefficient estimates equal to zero.

Thus regularization can simultaneously perform estimation and variable
selection.

A full treatment appears in statistical learning.

%%%%%%%%%%%%%%%%%%%%%%%%%%%%%%%%%%%%%%%%%%%%%%%%%%%%%%%%%%%%
\subsection{Robust Estimation}
\label{subsec:robust-estimation}
%%%%%%%%%%%%%%%%%%%%%%%%%%%%%%%%%%%%%%%%%%%%%%%%%%%%%%%%%%%%

An estimator is robust if its behavior does not change excessively under
small departures from assumptions or contamination by unusual
observations.

The sample mean is sensitive to extreme values.

The median is more resistant.

Thus robust estimation often trades some efficiency under an idealized
model for improved stability under model misspecification.

%%%%%%%%%%%%%%%%%%%%%%%%%%%%%%%%%%%%%%%%%%%%%%%%%%%%%%%%%%%%
\subsection{M-Estimators Preview}
\label{subsec:m-estimators-preview}
%%%%%%%%%%%%%%%%%%%%%%%%%%%%%%%%%%%%%%%%%%%%%%%%%%%%%%%%%%%%

An M-estimator minimizes an objective of the form

\[
\boxed{
\sum_{i=1}^{n}
\rho
\left(
\frac{
X_i-\theta
}{
s
}
\right),
}
\]

where \(\rho\) controls how residuals are penalized.

For squared loss,

\[
\rho(u)=u^2,
\]

the solution is the sample mean.

For absolute loss,

\[
\rho(u)=|u|,
\]

the solution is a median.

%%%%%%%%%%%%%%%%%%%%%%%%%%%%%%%%%%%%%%%%%%%%%%%%%%%%%%%%%%%%
\subsection{Estimating Equations}
\label{subsec:estimating-equations}
%%%%%%%%%%%%%%%%%%%%%%%%%%%%%%%%%%%%%%%%%%%%%%%%%%%%%%%%%%%%

Many estimators can be characterized as solutions of equations

\[
\boxed{
\sum_{i=1}^{n}
\psi(X_i,\theta)
=
0.
}
\]

Examples include:

\[
\boxed{
\text{score equations},
}
\]

and

\[
\boxed{
\text{M-estimation equations}.
}
\]

This provides a broad framework encompassing many classical and modern
estimators.

%%%%%%%%%%%%%%%%%%%%%%%%%%%%%%%%%%%%%%%%%%%%%%%%%%%%%%%%%%%%
\subsection{Asymptotic Efficiency}
\label{subsec:asymptotic-efficiency}
%%%%%%%%%%%%%%%%%%%%%%%%%%%%%%%%%%%%%%%%%%%%%%%%%%%%%%%%%%%%

Suppose two consistent estimators satisfy

\[
\sqrt{n}
(
T_{1n}-\theta
)
\xrightarrow{d}
N(0,V_1)
\]

and

\[
\sqrt{n}
(
T_{2n}-\theta
)
\xrightarrow{d}
N(0,V_2).
\]

One estimator is asymptotically more efficient if its asymptotic
variance is smaller.

A common relative efficiency is

\[
\boxed{
\frac{V_2}{V_1}.
}
\]

%%%%%%%%%%%%%%%%%%%%%%%%%%%%%%%%%%%%%%%%%%%%%%%%%%%%%%%%%%%%
\subsection{Asymptotic Efficiency of the MLE}
\label{subsec:mle-asymptotic-efficiency}
%%%%%%%%%%%%%%%%%%%%%%%%%%%%%%%%%%%%%%%%%%%%%%%%%%%%%%%%%%%%

Under regular parametric models, the MLE often has asymptotic variance

\[
\boxed{
\frac1{
nI_1(\theta)
}.
}
\]

This matches the asymptotic information lower bound.

Therefore regular MLEs are often asymptotically efficient.

This is one reason maximum likelihood is so widely used.

%%%%%%%%%%%%%%%%%%%%%%%%%%%%%%%%%%%%%%%%%%%%%%%%%%%%%%%%%%%%
\subsection{Delta Method for Estimators}
\label{subsec:delta-method-estimation}
%%%%%%%%%%%%%%%%%%%%%%%%%%%%%%%%%%%%%%%%%%%%%%%%%%%%%%%%%%%%

Suppose

\[
\sqrt{n}
(
\widehat{\theta}-\theta
)
\xrightarrow{d}
N(0,V).
\]

For differentiable \(g\),

\[
\boxed{
\sqrt{n}
[
g(\widehat{\theta})-g(\theta)
]
\xrightarrow{d}
N
\left(
0,
[g'(\theta)]^2V
\right).
}
\]

This allows standard errors for transformed parameters to be
approximated.

%%%%%%%%%%%%%%%%%%%%%%%%%%%%%%%%%%%%%%%%%%%%%%%%%%%%%%%%%%%%
\subsection{Worked Example: Estimated Odds}
\label{subsec:worked-delta-odds}
%%%%%%%%%%%%%%%%%%%%%%%%%%%%%%%%%%%%%%%%%%%%%%%%%%%%%%%%%%%%

\begin{workedexamplebox}{Odds from a Sample Proportion}

Suppose

\[
\widehat{p}
\]

estimates \(p\), and define the odds

\[
g(p)
=
\frac{p}{1-p}.
\]

Then

\[
g'(p)
=
\frac1{(1-p)^2}.
\]

If

\[
\operatorname{Var}(\widehat{p})
\approx
\frac{
p(1-p)
}{n},
\]

the delta method gives

\[
\operatorname{Var}
(
g(\widehat{p})
)
\approx
\left[
\frac1{(1-p)^2}
\right]^2
\frac{
p(1-p)
}{n}.
\]

Therefore,

\begin{answerbox}
\[
\boxed{
\operatorname{Var}
(
\widehat{\text{odds}}
)
\approx
\frac{
p
}{
n(1-p)^3
}.
}
\]
\end{answerbox}

\end{workedexamplebox}

%%%%%%%%%%%%%%%%%%%%%%%%%%%%%%%%%%%%%%%%%%%%%%%%%%%%%%%%%%%%
\subsection{Plug-In Principle}
\label{subsec:plugin-principle}
%%%%%%%%%%%%%%%%%%%%%%%%%%%%%%%%%%%%%%%%%%%%%%%%%%%%%%%%%%%%

A common estimation strategy replaces an unknown population quantity by
its empirical analogue.

For example,

\[
\mu
=
\mathbb{E}[X]
\]

is estimated by

\[
\overline{X}.
\]

A distributional functional

\[
T(F)
\]

may be estimated by

\[
\boxed{
T(F_n),
}
\]

where \(F_n\) is the empirical distribution.

This is called the plug-in principle.

%%%%%%%%%%%%%%%%%%%%%%%%%%%%%%%%%%%%%%%%%%%%%%%%%%%%%%%%%%%%
\subsection{Bootstrap Estimation}
\label{subsec:bootstrap-estimation}
%%%%%%%%%%%%%%%%%%%%%%%%%%%%%%%%%%%%%%%%%%%%%%%%%%%%%%%%%%%%

Bootstrap methods approximate estimator properties by repeatedly
resampling from the empirical distribution.

They can estimate:

\[
\boxed{
\text{standard errors},
}
\]

\[
\boxed{
\text{bias},
}
\]

and

\[
\boxed{
\text{sampling distributions}.
}
\]

A basic bootstrap bias estimate is

\[
\boxed{
\widehat{\operatorname{Bias}}_{\text{boot}}
=
\overline{\widehat{\theta}^*}
-
\widehat{\theta}.
}
\]

%%%%%%%%%%%%%%%%%%%%%%%%%%%%%%%%%%%%%%%%%%%%%%%%%%%%%%%%%%%%
\subsection{Jackknife Preview}
\label{subsec:jackknife-preview}
%%%%%%%%%%%%%%%%%%%%%%%%%%%%%%%%%%%%%%%%%%%%%%%%%%%%%%%%%%%%

The jackknife systematically removes one observation at a time.

Let

\[
\widehat{\theta}_{(-i)}
\]

be the estimate computed after deleting observation \(i\).

The collection

\[
\widehat{\theta}_{(-1)},
\ldots,
\widehat{\theta}_{(-n)}
\]

can be used to estimate bias and variance.

The jackknife predates the bootstrap and remains useful for smooth
statistics.

%%%%%%%%%%%%%%%%%%%%%%%%%%%%%%%%%%%%%%%%%%%%%%%%%%%%%%%%%%%%
\subsection{Estimator Sensitivity}
\label{subsec:estimator-sensitivity}
%%%%%%%%%%%%%%%%%%%%%%%%%%%%%%%%%%%%%%%%%%%%%%%%%%%%%%%%%%%%

An estimator may be sensitive to:

\[
\boxed{
\text{outliers},
}
\]

\[
\boxed{
\text{model misspecification},
}
\]

\[
\boxed{
\text{small sample size},
}
\]

\[
\boxed{
\text{parameter boundaries},
}
\]

and

\[
\boxed{
\text{dependence among observations}.
}
\]

Therefore numerical performance should be evaluated in addition to
formal optimality properties.

%%%%%%%%%%%%%%%%%%%%%%%%%%%%%%%%%%%%%%%%%%%%%%%%%%%%%%%%%%%%
\subsection{Model Misspecification}
\label{subsec:model-misspecification-estimation}
%%%%%%%%%%%%%%%%%%%%%%%%%%%%%%%%%%%%%%%%%%%%%%%%%%%%%%%%%%%%

Maximum likelihood assumes the chosen parametric family adequately
describes the data-generating mechanism.

If the true distribution is outside the model family, the MLE may
converge to a parameter value that best approximates the true
distribution in a specific information-theoretic sense rather than to a
literal true parameter.

Thus a mathematically optimal estimator under the wrong model can still
be scientifically misleading.

%%%%%%%%%%%%%%%%%%%%%%%%%%%%%%%%%%%%%%%%%%%%%%%%%%%%%%%%%%%%
\subsection{Estimating Under Dependence}
\label{subsec:estimation-under-dependence}
%%%%%%%%%%%%%%%%%%%%%%%%%%%%%%%%%%%%%%%%%%%%%%%%%%%%%%%%%%%%

Many formulas in this section assume independent observations.

For dependent data,

\[
\operatorname{Var}
\left(
\frac1n
\sum X_i
\right)
\]

includes covariance terms:

\[
\boxed{
\frac1{n^2}
\left[
\sum_i
\operatorname{Var}(X_i)
+
2
\sum_{i<j}
\operatorname{Cov}(X_i,X_j)
\right].
}
\]

Time series, clustered samples, and repeated measurements require
appropriate dependence-aware methods.

%%%%%%%%%%%%%%%%%%%%%%%%%%%%%%%%%%%%%%%%%%%%%%%%%%%%%%%%%%%%
\subsection{Estimation Workflow}
\label{subsec:estimation-workflow}
%%%%%%%%%%%%%%%%%%%%%%%%%%%%%%%%%%%%%%%%%%%%%%%%%%%%%%%%%%%%

A systematic estimation workflow is:

\begin{enumerate}
    \item define the parameter of interest;
    \item specify the probability model;
    \item check parameter constraints;
    \item choose an estimation principle;
    \item derive the estimator;
    \item examine bias;
    \item compute or approximate variance;
    \item evaluate mean squared error;
    \item check consistency;
    \item examine efficiency when relevant;
    \item identify sufficient statistics when possible;
    \item compute standard errors;
    \item inspect robustness to unusual observations;
    \item verify model assumptions;
    \item connect the point estimate to interval estimation or testing.
\end{enumerate}

%%%%%%%%%%%%%%%%%%%%%%%%%%%%%%%%%%%%%%%%%%%%%%%%%%%%%%%%%%%%
\subsection{Common Errors}
\label{subsec:statistical-estimation-common-errors}
%%%%%%%%%%%%%%%%%%%%%%%%%%%%%%%%%%%%%%%%%%%%%%%%%%%%%%%%%%%%

\subsubsection{Confusing Parameter and Estimate}

A parameter is a fixed unknown population quantity.

An estimate is a numerical value calculated from observed data.

%%%%%%%%%%%%%%%%%%%%%%%%%%%%%%%%%%%%%%%%%%%%%%%%%%%%%%%%%%%%
\subsubsection{Confusing Estimator and Estimate}

The estimator is random before sampling.

The estimate is the realized numerical value.

%%%%%%%%%%%%%%%%%%%%%%%%%%%%%%%%%%%%%%%%%%%%%%%%%%%%%%%%%%%%
\subsubsection{Assuming Unbiased Means Best}

An unbiased estimator can have very large variance.

Mean squared error considers both bias and variance.

%%%%%%%%%%%%%%%%%%%%%%%%%%%%%%%%%%%%%%%%%%%%%%%%%%%%%%%%%%%%
\subsubsection{Assuming a Biased Estimator Is Bad}

A biased estimator may have substantially smaller MSE.

Shrinkage estimators provide important examples.

%%%%%%%%%%%%%%%%%%%%%%%%%%%%%%%%%%%%%%%%%%%%%%%%%%%%%%%%%%%%
\subsubsection{Confusing Consistency with Unbiasedness}

Consistency is an asymptotic property.

Unbiasedness is an expectation property at a given sample size.

%%%%%%%%%%%%%%%%%%%%%%%%%%%%%%%%%%%%%%%%%%%%%%%%%%%%%%%%%%%%
\subsubsection{Ignoring Parameter Constraints}

An estimator must belong to the parameter space.

For example,

\[
p
\]

must satisfy

\[
0\leq p\leq1.
\]

%%%%%%%%%%%%%%%%%%%%%%%%%%%%%%%%%%%%%%%%%%%%%%%%%%%%%%%%%%%%
\subsubsection{Treating Likelihood as \(P(\theta\mid X)\)}

In frequentist inference, likelihood is a function of the parameter for
fixed observed data.

It is not a posterior probability distribution.

%%%%%%%%%%%%%%%%%%%%%%%%%%%%%%%%%%%%%%%%%%%%%%%%%%%%%%%%%%%%
\subsubsection{Setting the Score to Zero Without Checking the Maximum}

A score equation identifies candidate extrema.

One must still check:

\[
\boxed{
\text{second derivatives},
}
\]

\[
\boxed{
\text{boundaries},
}
\]

or

\[
\boxed{
\text{direct likelihood behavior}.
}
\]

%%%%%%%%%%%%%%%%%%%%%%%%%%%%%%%%%%%%%%%%%%%%%%%%%%%%%%%%%%%%
\subsubsection{Ignoring Support That Depends on the Parameter}

In models such as

\[
\operatorname{Uniform}(0,\theta),
\]

the support depends on \(\theta\).

Standard differentiation-based regularity arguments may fail.

%%%%%%%%%%%%%%%%%%%%%%%%%%%%%%%%%%%%%%%%%%%%%%%%%%%%%%%%%%%%
\subsubsection{Assuming Every MLE Is Unbiased}

MLEs are often biased at finite sample sizes.

Their major strengths are often consistency and asymptotic efficiency.

%%%%%%%%%%%%%%%%%%%%%%%%%%%%%%%%%%%%%%%%%%%%%%%%%%%%%%%%%%%%
\subsubsection{Using the Cram\'er--Rao Bound Without Checking Conditions}

The classical Cram\'er--Rao bound requires regularity conditions.

It may not apply directly to irregular models or parameter-dependent
support.

%%%%%%%%%%%%%%%%%%%%%%%%%%%%%%%%%%%%%%%%%%%%%%%%%%%%%%%%%%%%
\subsubsection{Confusing Observed and Expected Information}

Observed information depends on the realized sample.

Fisher information averages over the sampling model.

%%%%%%%%%%%%%%%%%%%%%%%%%%%%%%%%%%%%%%%%%%%%%%%%%%%%%%%%%%%%
\subsubsection{Ignoring Nuisance Parameters}

Uncertainty in nuisance parameters can affect the uncertainty of the
parameter of interest.

%%%%%%%%%%%%%%%%%%%%%%%%%%%%%%%%%%%%%%%%%%%%%%%%%%%%%%%%%%%%
\subsubsection{Assuming More Complex Means Better}

An estimator should be judged by its statistical properties,
interpretability, robustness, and relevance to the scientific problem.

%%%%%%%%%%%%%%%%%%%%%%%%%%%%%%%%%%%%%%%%%%%%%%%%%%%%%%%%%%%%
\subsection{Applications}
\label{subsec:statistical-estimation-applications}
%%%%%%%%%%%%%%%%%%%%%%%%%%%%%%%%%%%%%%%%%%%%%%%%%%%%%%%%%%%%

\begin{applicationbox}

\textbf{Survey Research.}
Population means, totals, proportions, and subgroup quantities are
estimated from samples.

\medskip

\textbf{Scientific Experiments.}
Treatment effects, response rates, and model parameters are estimated
from experimental observations.

\medskip

\textbf{Engineering.}
Failure rates, process means, variances, and reliability parameters are
estimated from observed systems.

\medskip

\textbf{Environmental Science.}
Researchers estimate concentrations, trends, spatial effects, and model
parameters from environmental measurements.

\medskip

\textbf{Biology and Epidemiology.}
Population growth rates, disease risks, survival parameters, and
treatment effects require statistical estimation.

\medskip

\textbf{Economics and Finance.}
Regression coefficients, volatility, returns, and structural parameters
are estimated from observational data.

\medskip

\textbf{Machine Learning.}
Training a statistical model usually means estimating parameters from
data, often through likelihood or regularized loss minimization.

\medskip

\textbf{Bayesian Statistics.}
Parameters are estimated through posterior distributions combining prior
information and likelihood.

\end{applicationbox}

%%%%%%%%%%%%%%%%%%%%%%%%%%%%%%%%%%%%%%%%%%%%%%%%%%%%%%%%%%%%
\subsection{Exercises}
\label{subsec:statistical-estimation-exercises}
%%%%%%%%%%%%%%%%%%%%%%%%%%%%%%%%%%%%%%%%%%%%%%%%%%%%%%%%%%%%

\begin{exercise}[1]
Define an estimator.
\end{exercise}

\begin{exercise}[1]
Distinguish an estimator from an estimate.
\end{exercise}

\begin{exercise}[1]
Define bias.
\end{exercise}

\begin{exercise}[1]
Define consistency.
\end{exercise}

\begin{exercise}[1]
Define mean squared error.
\end{exercise}

\begin{exercise}[1]
Define a sufficient statistic.
\end{exercise}

\begin{exercise}[1]
State the factorization theorem.
\end{exercise}

\begin{exercise}[1]
Define the likelihood function.
\end{exercise}

\begin{exercise}[1]
Define the maximum likelihood estimator.
\end{exercise}

\begin{exercise}[1]
Define the score function.
\end{exercise}

\begin{exercise}[1]
Define Fisher information.
\end{exercise}

\begin{exercise}[1]
State the Cram\'er--Rao lower bound.
\end{exercise}

\begin{exercise}[2]
Show that the sample mean is unbiased for the population mean.
\end{exercise}

\begin{exercise}[2]
Show that

\[
\frac1n
\sum_{i=1}^{n}
(X_i-\overline{X})^2
\]

is biased for \(\sigma^2\).
\end{exercise}

\begin{exercise}[2]
Suppose an estimator has bias \(2\) and variance \(9\). Find its MSE.
\end{exercise}

\begin{exercise}[2]
Suppose estimator \(T_1\) has variance \(4\) and estimator \(T_2\) has
variance \(10\), and both are unbiased. Which is more efficient?
\end{exercise}

\begin{exercise}[2]
Find the method-of-moments estimator of \(\lambda\) for

\[
X_i\sim\operatorname{Poisson}(\lambda).
\]
\end{exercise}

\begin{exercise}[2]
Find the method-of-moments estimator of \(\theta\) for

\[
X_i\sim\operatorname{Uniform}(0,\theta).
\]
\end{exercise}

\begin{exercise}[2]
Find the MLE of \(p\) for a Bernoulli sample.
\end{exercise}

\begin{exercise}[2]
Find the MLE of \(\lambda\) for a Poisson sample.
\end{exercise}

\begin{exercise}[2]
Find the MLE of the exponential rate parameter.
\end{exercise}

\begin{exercise}[2]
Find the MLE of a normal mean when variance is known.
\end{exercise}

\begin{exercise}[2]
Find the MLE of \(\theta\) for

\[
X_i\sim\operatorname{Uniform}(0,\theta).
\]
\end{exercise}

\begin{exercise}[3]
Derive the bias--variance decomposition of MSE.
\end{exercise}

\begin{exercise}[3]
Show that if

\[
\operatorname{Bias}(\widehat{\theta}_n)\to0
\]

and

\[
\operatorname{Var}(\widehat{\theta}_n)\to0,
\]

then \(\widehat{\theta}_n\) is consistent.
\end{exercise}

\begin{exercise}[3]
Use the factorization theorem to show that

\[
\sum X_i
\]

is sufficient for \(p\) in a Bernoulli model.
\end{exercise}

\begin{exercise}[3]
Use the factorization theorem to identify a sufficient statistic for
\(\lambda\) in a Poisson model.
\end{exercise}

\begin{exercise}[3]
Show that \(\overline{X}\) is sufficient for \(\mu\) in a normal model
with known variance.
\end{exercise}

\begin{exercise}[3]
Derive the MLEs of both \(\mu\) and \(\sigma^2\) for a normal sample.
\end{exercise}

\begin{exercise}[3]
Show that the normal variance MLE is biased.
\end{exercise}

\begin{exercise}[3]
Show that the normal variance MLE is consistent.
\end{exercise}

\begin{exercise}[3]
Derive the score function for Bernoulli data.
\end{exercise}

\begin{exercise}[3]
Derive the Fisher information for Bernoulli data.
\end{exercise}

\begin{exercise}[3]
Derive the Fisher information for Poisson data.
\end{exercise}

\begin{exercise}[3]
Derive the Fisher information for a normal mean with known variance.
\end{exercise}

\begin{exercise}[3]
Show that the Bernoulli sample proportion attains the Cram\'er--Rao lower
bound.
\end{exercise}

\begin{exercise}[3]
For

\[
X_i\sim N(\mu,\sigma^2)
\]

with known \(\sigma^2\), show that \(\overline{X}\) attains the
Cram\'er--Rao bound for estimating \(\mu\).
\end{exercise}

\begin{exercise}[3]
Use the invariance property of MLEs to find the MLE of

\[
e^\mu
\]

if \(\widehat{\mu}\) is the MLE of \(\mu\).
\end{exercise}

\begin{exercise}[3]
Suppose

\[
\widehat{\theta}
\]

is asymptotically normal. Apply the delta method to

\[
g(\theta)=1/\theta.
\]
\end{exercise}

\begin{exercise}[3]
Explain why an MLE may occur at the boundary of the parameter space.
\end{exercise}

\begin{exercise}[3]
Explain why the regular Cram\'er--Rao argument is problematic for
\(\operatorname{Uniform}(0,\theta)\).
\end{exercise}

\begin{exercise}[3]
Show using the law of total variance why Rao--Blackwellization cannot
increase variance for an unbiased estimator.
\end{exercise}

\begin{exercise}[3]
Explain how a complete sufficient statistic leads to uniqueness of an
UMVU estimator.
\end{exercise}

\begin{exercise}[3]
Construct an example where a biased estimator has smaller MSE than an
unbiased estimator.
\end{exercise}

\begin{exercise}[3]
Compare MLE and method-of-moments estimators for the upper endpoint of a
uniform distribution.
\end{exercise}

\begin{exercise}[3]
Explain the difference between expected Fisher information and observed
information.
\end{exercise}

\begin{exercise}[4]
Prove the Cram\'er--Rao lower bound under standard regularity conditions.
\end{exercise}

\begin{exercise}[4]
Study minimal sufficient statistics using likelihood ratios.
\end{exercise}

\begin{exercise}[4]
Investigate complete sufficient statistics for one-parameter exponential
families.
\end{exercise}

\begin{exercise}[4]
Apply the Lehmann--Scheff\'e theorem to construct UMVU estimators.
\end{exercise}

\begin{exercise}[4]
Study the asymptotic consistency of maximum likelihood estimators.
\end{exercise}

\begin{exercise}[4]
Derive the asymptotic normality of MLEs using a Taylor expansion of the
score function.
\end{exercise}

\begin{exercise}[4]
Study the multiparameter Cram\'er--Rao inequality.
\end{exercise}

\begin{exercise}[4]
Investigate profile likelihood and nuisance parameters.
\end{exercise}

\begin{exercise}[4]
Study estimating equations and M-estimators.
\end{exercise}

\begin{exercise}[4]
Investigate Huber robust estimation.
\end{exercise}

\begin{exercise}[4]
Study Bayesian estimators under squared, absolute, and zero-one loss.
\end{exercise}

\begin{exercise}[4]
Investigate shrinkage estimation and the James--Stein phenomenon.
\end{exercise}

\begin{exercise}[4]
Study bootstrap bias correction.
\end{exercise}

\begin{exercise}[4]
Investigate jackknife variance estimation.
\end{exercise}

\begin{exercise}[4]
Study estimator behavior under model misspecification.
\end{exercise}

%%%%%%%%%%%%%%%%%%%%%%%%%%%%%%%%%%%%%%%%%%%%%%%%%%%%%%%%%%%%
\subsection{Conceptual Summary}
\label{subsec:statistical-estimation-summary}
%%%%%%%%%%%%%%%%%%%%%%%%%%%%%%%%%%%%%%%%%%%%%%%%%%%%%%%%%%%%

An estimator is a statistic

\[
\boxed{
\widehat{\theta}
=
T(X_1,\ldots,X_n)
}
\]

used to estimate an unknown parameter.

Its bias is

\[
\boxed{
\operatorname{Bias}(\widehat{\theta})
=
\mathbb{E}[\widehat{\theta}]
-
\theta.
}
\]

Its mean squared error is

\[
\boxed{
\operatorname{MSE}(\widehat{\theta})
=
\operatorname{Var}(\widehat{\theta})
+
\operatorname{Bias}(\widehat{\theta})^2.
}
\]

Consistency means

\[
\boxed{
\widehat{\theta}_n
\xrightarrow{P}
\theta.
}
\]

A sufficient statistic compresses the sample without losing
parameter-relevant information.

The factorization theorem identifies sufficiency through

\[
\boxed{
f(x;\theta)
=
g(T(x),\theta)h(x).
}
\]

The maximum likelihood estimator satisfies

\[
\boxed{
\widehat{\theta}_{MLE}
=
\arg\max_{\theta}
L(\theta).
}
\]

The score is

\[
\boxed{
U(\theta)
=
\ell'(\theta).
}
\]

Fisher information is

\[
\boxed{
I(\theta)
=
\mathbb{E}[U(\theta)^2]
=
-
\mathbb{E}[\ell''(\theta)]
}
\]

under regularity conditions.

The Cram\'er--Rao bound is

\[
\boxed{
\operatorname{Var}(\widehat{\theta})
\geq
\frac1{I_n(\theta)}
}
\]

for regular unbiased estimation of a scalar parameter.

Regular MLEs often satisfy

\[
\boxed{
\sqrt{n}
(
\widehat{\theta}_{MLE}-\theta
)
\xrightarrow{d}
N
\left(
0,
\frac1{I_1(\theta)}
\right).
}
\]

Statistical estimation therefore studies how sample data are converted
into parameter estimates and how the quality of those estimates is
measured mathematically.

%%%%%%%%%%%%%%%%%%%%%%%%%%%%%%%%%%%%%%%%%%%%%%%%%%%%%%%%%%%%
\subsection{Section Connections}
\label{subsec:statistical-estimation-connections}
%%%%%%%%%%%%%%%%%%%%%%%%%%%%%%%%%%%%%%%%%%%%%%%%%%%%%%%%%%%%

\chapterconnection{
Statistical Estimation builds directly on the sampling distributions of
Section~\ref{sec:sampling-distributions}. Sampling distributions provide
the bias, variance, standard errors, and asymptotic behavior used to
evaluate estimators. Likelihood connects probability models with
parameter estimation, while sufficient statistics, Fisher information,
and efficiency formalize how much information a sample contains about an
unknown parameter. The next section, Confidence Intervals, uses point
estimators and their sampling distributions to construct ranges of
plausible parameter values. Hypothesis Testing then uses the same
likelihood, information, and sampling-distribution ideas to evaluate
statistical claims.
}

%%%%%%%%%%%%%%%%%%%%%%%%%%%%%%%%%%%%%%%%%%%%%%%%%%%%%%%%%%%%
% SECTION SOURCE TRACKING
%%%%%%%%%%%%%%%%%%%%%%%%%%%%%%%%%%%%%%%%%%%%%%%%%%%%%%%%%%%%

%------------------------------------------------------------
% 8.4 Statistical Estimation
%------------------------------------------------------------
%
% Source basis:
%   - Standard mathematical statistics.
%   - Standard likelihood theory.
%   - Standard decision-theoretic and asymptotic estimation concepts.
%
% Used for:
%   - point estimation;
%   - parameters, estimators, and estimates;
%   - bias;
%   - variance;
%   - mean squared error;
%   - bias--variance decomposition;
%   - consistency;
%   - mean-square consistency;
%   - efficiency;
%   - sufficient statistics;
%   - factorization theorem;
%   - minimal sufficiency preview;
%   - completeness;
%   - Rao--Blackwell theorem;
%   - Lehmann--Scheffé theorem;
%   - UMVU estimators;
%   - method of moments;
%   - likelihood and log-likelihood;
%   - maximum likelihood estimation;
%   - Bernoulli, Poisson, exponential, normal, and uniform MLEs;
%   - boundary likelihood estimates;
%   - invariance of MLEs;
%   - score functions;
%   - Fisher information;
%   - Cramér--Rao lower bound;
%   - estimator efficiency;
%   - observed information;
%   - MLE consistency;
%   - asymptotic normality;
%   - likelihood curvature;
%   - relative and profile likelihoods;
%   - nuisance parameters;
%   - multiparameter estimation;
%   - Fisher information matrices;
%   - identifiability;
%   - constrained estimation;
%   - loss and risk;
%   - Bayesian estimation preview;
%   - MAP estimation;
%   - shrinkage and regularization;
%   - robust estimation;
%   - M-estimators;
%   - estimating equations;
%   - delta-method estimation;
%   - plug-in principle;
%   - bootstrap and jackknife estimation;
%   - model misspecification;
%   - estimation under dependence;
%   - applications and exercises.
%
% Notes:
%   - Confidence Intervals are developed next in Section 8.5.
%   - Hypothesis Testing is developed in Section 8.6.
%   - Mathematical Statistics in Section 8.13 develops sufficiency,
%     completeness, information, decision theory, and asymptotics more
%     deeply.
%   - Bayesian estimation is developed fully in Section 8.14.
%
%%%%%%%%%%%%%%%%%%%%%%%%%%%%%%%%%%%%%%%%%%%%%%%%%%%%%%%%%%%%